\section{Part A - Linux Commands
Explanation}\label{part-a---linux-commands-explanation}

\subsection{What will the following commands
do?}\label{what-will-the-following-commands-do}

\subsubsection{\texorpdfstring{\texttt{echo\ "Hello,\ World!"}}{echo "Hello, World!"}}\label{echo-hello-world}

Prints \texttt{Hello,\ World!} to the terminal.

\subsubsection{\texorpdfstring{\texttt{name="Productive"}}{name="Productive"}}\label{nameproductive}

Creates a variable \texttt{name} and assigns it the value
\texttt{Productive}.

\subsubsection{\texorpdfstring{\texttt{touch\ file.txt}}{touch file.txt}}\label{touch-file.txt}

Creates an empty file named \texttt{file.txt} or updates its timestamp
if it already exists.

\subsubsection{\texorpdfstring{\texttt{ls\ -a}}{ls -a}}\label{ls--a}

Lists all files and directories in the current directory, including
hidden ones (those starting with \texttt{.}).

\subsubsection{\texorpdfstring{\texttt{rm\ file.txt}}{rm file.txt}}\label{rm-file.txt}

Removes the file \texttt{file.txt} permanently.

\subsubsection{\texorpdfstring{\texttt{cp\ file1.txt\ file2.txt}}{cp file1.txt file2.txt}}\label{cp-file1.txt-file2.txt}

Copies \texttt{file1.txt} to \texttt{file2.txt}. If \texttt{file2.txt}
exists, it will be overwritten.

\subsubsection{\texorpdfstring{\texttt{mv\ file.txt\ /path/to/directory/}}{mv file.txt /path/to/directory/}}\label{mv-file.txt-pathtodirectory}

Moves \texttt{file.txt} to the specified directory.

\subsubsection{\texorpdfstring{\texttt{chmod\ 755\ script.sh}}{chmod 755 script.sh}}\label{chmod-755-script.sh}

Grants the owner full permissions (read, write, execute) and gives
others read and execute permissions on \texttt{script.sh}.

\subsubsection{\texorpdfstring{\texttt{grep\ "pattern"\ file.txt}}{grep "pattern" file.txt}}\label{grep-pattern-file.txt}

Searches for occurrences of ``pattern'' in \texttt{file.txt} and prints
matching lines.

\subsubsection{\texorpdfstring{\texttt{kill\ PID}}{kill PID}}\label{kill-pid}

Terminates the process with the specified Process ID (PID).

\subsubsection{\texorpdfstring{\texttt{mkdir\ mydir\ \&\&\ cd\ mydir\ \&\&\ touch\ file.txt\ \&\&\ echo\ "Hello,\ World!"\ \textgreater{}\ file.txt\ \&\&\ cat\ file.txt}}{mkdir mydir \&\& cd mydir \&\& touch file.txt \&\& echo "Hello, World!" \textgreater{} file.txt \&\& cat file.txt}}\label{mkdir-mydir-cd-mydir-touch-file.txt-echo-hello-world-file.txt-cat-file.txt}

\begin{itemize}
\tightlist
\item
  Creates a directory \texttt{mydir}
\item
  Changes into \texttt{mydir}
\item
  Creates an empty file \texttt{file.txt}
\item
  Writes ``Hello, World!'' into \texttt{file.txt}
\item
  Displays the contents of \texttt{file.txt}
\end{itemize}

\subsubsection{\texorpdfstring{\texttt{ls\ -l\ \textbar{}\ grep\ ".txt"}}{ls -l \textbar{} grep ".txt"}}\label{ls--l-grep-.txt}

Lists files in long format and filters only those containing ``. Txt''
in their names.

\subsubsection{\texorpdfstring{\texttt{cat\ file1.txt\ file2.txt\ \textbar{}\ sort\ \textbar{}\ uniq}}{cat file1.txt file2.txt \textbar{} sort \textbar{} uniq}}\label{cat-file1.txt-file2.txt-sort-uniq}

Concatenates \texttt{file1.txt} and \texttt{file2.txt}, sorts them, and
removes duplicate lines.

\subsubsection{\texorpdfstring{\texttt{ls\ -l\ \textbar{}\ grep\ "\^{}d"}}{ls -l \textbar{} grep "\^{}d"}}\label{ls--l-grep-d}

Lists directories (entries starting with \texttt{d} in long format
output).

\subsubsection{\texorpdfstring{\texttt{grep\ -r\ "pattern"\ /path/to/directory/}}{grep -r "pattern" /path/to/directory/}}\label{grep--r-pattern-pathtodirectory}

Searches for ``pattern'' recursively in all files under
\texttt{/path/to/directory/}.

\subsubsection{\texorpdfstring{\texttt{cat\ file1.txt\ file2.txt\ \textbar{}\ sort\ \textbar{}\ uniq\ -d}}{cat file1.txt file2.txt \textbar{} sort \textbar{} uniq -d}}\label{cat-file1.txt-file2.txt-sort-uniq--d}

Concatenates \texttt{file1.txt} and \texttt{file2.txt}, sorts them, and
displays only duplicate lines.

\subsubsection{\texorpdfstring{\texttt{chmod\ 644\ file.txt}}{chmod 644 file.txt}}\label{chmod-644-file.txt}

Grants the owner read and write permissions, while others get read-only
access to \texttt{file.txt}.

\subsubsection{\texorpdfstring{\texttt{cp\ -r\ source\_directory\ destination\_directory}}{cp -r source\_directory destination\_directory}}\label{cp--r-source_directory-destination_directory}

Recursively copies \texttt{source\_directory} to
\texttt{destination\_directory}, preserving contents.

\subsubsection{\texorpdfstring{\texttt{find\ /path/to/search\ -name\ "*.txt"}}{find /path/to/search -name "*.txt"}}\label{find-pathtosearch--name-.txt}

Finds all \texttt{.txt} files in \texttt{/path/to/search} and its
subdirectories.

\subsubsection{\texorpdfstring{\texttt{chmod\ u+x\ file.txt}}{chmod u+x file.txt}}\label{chmod-ux-file.txt}

Gives the owner (\texttt{u}) execute permission on \texttt{file.txt}.

\subsubsection{\texorpdfstring{\texttt{echo\ \$PATH}}{echo \$PATH}}\label{echo-path}

Displays the system's \texttt{PATH} environment variable, listing
directories where executable files are searched for.

\begin{center}\rule{0.5\linewidth}{0.5pt}\end{center}

\section{Part B - Identify True or
False}\label{part-b---identify-true-or-false}

\begin{enumerate}
\def\labelenumi{\arabic{enumi}.}
\tightlist
\item
  \textbf{True} - \texttt{ls} is used to list files and directories in a
  directory.
\item
  \textbf{True} - \texttt{mv} is used to move files and directories.
\item
  \textbf{False} - \texttt{cd} is used to change directories, not copy
  files and directories.
\item
  \textbf{True} - \texttt{pwd} stands for ``print working directory''
  and displays the current directory.
\item
  \textbf{True} - \texttt{grep} is used to search for patterns in files.
\item
  \textbf{True} - \texttt{chmod\ 755\ file.txt} gives read, write, and
  execute permissions to the owner, and read and execute permissions to
  group and others.
\item
  \textbf{True} - \texttt{mkdir\ -p\ directory1/directory2} creates
  nested directories, creating \texttt{directory2} inside
  \texttt{directory1} if \texttt{directory1} does not exist.
\item
  \textbf{True} - \texttt{rm\ -rf\ file.txt} deletes a file forcefully
  without confirmation.
\end{enumerate}

\begin{longtable}[]{@{}
  >{\raggedright\arraybackslash}p{(\linewidth - 0\tabcolsep) * \real{0.0556}}@{}}
\toprule\noalign{}
\endhead
\bottomrule\noalign{}
\endlastfoot
1. \textbf{Incorrect} - \texttt{chmodx} is not a valid command. The
correct command to change file permissions is \texttt{chmod}. 2.
\textbf{Incorrect} - \texttt{cpy} is not a valid command. The correct
command to copy files and directories is \texttt{cp}. 3.
\textbf{Incorrect} - \texttt{mkfile} is not a standard Linux command. To
create a new file, use \texttt{touch\ filename}. 4. \textbf{Incorrect} -
\texttt{catx} is not a valid command. The correct command to concatenate
files is \texttt{cat}. 5. \textbf{Incorrect} - \texttt{rn} is not a
valid command. To rename files, use the \texttt{mv} command
(\texttt{mv\ oldname\ newname}). \\
\end{longtable}

\section{Part C - Shell Scripting
Questions}\label{part-c---shell-scripting-questions}

\subsubsection{Question 1: Write a shell script that prints ``Hello,
World!'' to the
terminal.}\label{question-1-write-a-shell-script-that-prints-hello-world-to-the-terminal.}

\begin{Shaded}
\begin{Highlighting}[]
\CommentTok{\#!/bin/bash}
\BuiltInTok{echo} \StringTok{"Hello, World!"}
\end{Highlighting}
\end{Shaded}

!{[}{[}Pasted image 20250301213434.png{]}{]} \#\#\# Question 2: Declare
a variable named ``name'' and assign the value ``CDAC Mumbai'' to it.
Print the value of the variable.

\begin{Shaded}
\begin{Highlighting}[]
\CommentTok{\#!/bin/bash}
\VariableTok{name}\OperatorTok{=}\StringTok{"CDAC Mumbai"}
\BuiltInTok{echo} \StringTok{"}\VariableTok{$name}\StringTok{"}
\end{Highlighting}
\end{Shaded}

!{[}{[}Pasted image 20250301213539.png{]}{]} \#\#\# Question 3: Write a
shell script that takes a number as input from the user and prints it.

\begin{Shaded}
\begin{Highlighting}[]
\CommentTok{\#!/bin/bash}
\BuiltInTok{read} \AttributeTok{{-}p} \StringTok{"Enter a number: "} \VariableTok{num}
\BuiltInTok{echo} \StringTok{"You entered: }\VariableTok{$num}\StringTok{"}
\end{Highlighting}
\end{Shaded}

!{[}{[}Pasted image 20250301214517.png{]}{]} \#\#\# Question 4: Write a
shell script that performs addition of two numbers (e.g., 5 and 3) and
prints the result.

\begin{Shaded}
\begin{Highlighting}[]
\CommentTok{\#!/bin/bash}
\VariableTok{a}\OperatorTok{=}\NormalTok{5}
\VariableTok{b}\OperatorTok{=}\NormalTok{3}
\VariableTok{sum}\OperatorTok{=}\VariableTok{$((a} \OperatorTok{+} \VariableTok{b))}
\BuiltInTok{echo} \StringTok{"Sum: }\VariableTok{$sum}\StringTok{"}
\end{Highlighting}
\end{Shaded}

!{[}{[}Pasted image 20250301215138.png{]}{]}

\subsubsection{Question 5: Write a shell script that takes a number as
input and prints ``Even'' if it is even, otherwise prints
``Odd''.}\label{question-5-write-a-shell-script-that-takes-a-number-as-input-and-prints-even-if-it-is-even-otherwise-prints-odd.}

\begin{Shaded}
\begin{Highlighting}[]
\CommentTok{\#!/bin/bash}
\BuiltInTok{read} \AttributeTok{{-}p} \StringTok{"Enter a number: "} \VariableTok{num}
\ControlFlowTok{if} \KeywordTok{((}\VariableTok{num} \OperatorTok{\%} \DecValTok{2} \OperatorTok{==} \DecValTok{0}\KeywordTok{));} \ControlFlowTok{then}
    \BuiltInTok{echo} \StringTok{"Even"}
\ControlFlowTok{else}
    \BuiltInTok{echo} \StringTok{"Odd"}
\ControlFlowTok{fi}
\end{Highlighting}
\end{Shaded}

!{[}{[}Pasted image 20250301215326.png{]}{]} \#\#\# Question 6: Write a
shell script that uses a for loop to print numbers from 1 to 5.

\begin{Shaded}
\begin{Highlighting}[]
\CommentTok{\#!/bin/bash}
\ControlFlowTok{for}\NormalTok{ i }\KeywordTok{in} \DataTypeTok{\{}\DecValTok{1}\DataTypeTok{..}\DecValTok{5}\DataTypeTok{\}}\KeywordTok{;} \ControlFlowTok{do}
    \BuiltInTok{echo} \StringTok{"}\VariableTok{$i}\StringTok{"}
\ControlFlowTok{done}
\end{Highlighting}
\end{Shaded}

!{[}{[}Pasted image 20250301215226.png{]}{]} \#\#\# Question 7: Write a
shell script that uses a while loop to print numbers from 1 to 5.

\begin{Shaded}
\begin{Highlighting}[]
\CommentTok{\#!/bin/bash}
\VariableTok{i}\OperatorTok{=}\NormalTok{1}
\ControlFlowTok{while} \BuiltInTok{[} \VariableTok{$i} \OtherTok{{-}le}\NormalTok{ 5 }\BuiltInTok{]}\KeywordTok{;} \ControlFlowTok{do}
    \BuiltInTok{echo} \StringTok{"}\VariableTok{$i}\StringTok{"}
    \KeywordTok{((}\VariableTok{i}\OperatorTok{++}\KeywordTok{))}
\ControlFlowTok{done}
\end{Highlighting}
\end{Shaded}

!{[}{[}Pasted image 20250301215433.png{]}{]} \#\#\# Question 8: Write a
shell script that checks if a file named ``file. Txt'' exists in the
current directory.

\begin{Shaded}
\begin{Highlighting}[]
\CommentTok{\#!/bin/bash}
\ControlFlowTok{if} \BuiltInTok{[} \OtherTok{{-}f} \StringTok{"file.txt"} \BuiltInTok{]}\KeywordTok{;} \ControlFlowTok{then}
    \BuiltInTok{echo} \StringTok{"File exists"}
\ControlFlowTok{else}
    \BuiltInTok{echo} \StringTok{"File does not exist"}
\ControlFlowTok{fi}
\end{Highlighting}
\end{Shaded}

!{[}{[}Pasted image 20250301215713.png{]}{]} \#\#\# Question 9: Write a
shell script that checks if a number is greater than 10 and prints a
message accordingly.

\begin{Shaded}
\begin{Highlighting}[]
\CommentTok{\#!/bin/bash}
\BuiltInTok{read} \AttributeTok{{-}p} \StringTok{"Enter a number: "} \VariableTok{num}
\ControlFlowTok{if} \BuiltInTok{[} \VariableTok{$num} \OtherTok{{-}gt}\NormalTok{ 10 }\BuiltInTok{]}\KeywordTok{;} \ControlFlowTok{then}
    \BuiltInTok{echo} \StringTok{"Number is greater than 10"}
\ControlFlowTok{else}
    \BuiltInTok{echo} \StringTok{"Number is 10 or less"}
\ControlFlowTok{fi}
\end{Highlighting}
\end{Shaded}

!{[}{[}Pasted image 20250301215936.png{]}{]} \#\#\# Question 10: Write a
shell script that prints a multiplication table for numbers from 1 to 5.

\begin{Shaded}
\begin{Highlighting}[]
\CommentTok{\#!/bin/bash}
\ControlFlowTok{for}\NormalTok{ i }\KeywordTok{in} \DataTypeTok{\{}\DecValTok{1}\DataTypeTok{..}\DecValTok{5}\DataTypeTok{\}}\KeywordTok{;} \ControlFlowTok{do}
    \ControlFlowTok{for}\NormalTok{ j }\KeywordTok{in} \DataTypeTok{\{}\DecValTok{1}\DataTypeTok{..}\DecValTok{5}\DataTypeTok{\}}\KeywordTok{;} \ControlFlowTok{do}
        \BuiltInTok{printf} \StringTok{"\%4d"} \VariableTok{$((i} \OperatorTok{*} \VariableTok{j))}
    \ControlFlowTok{done}
    \BuiltInTok{echo}
\ControlFlowTok{done}
\end{Highlighting}
\end{Shaded}

!{[}{[}Pasted image 20250301220214.png{]}{]} !{[}{[}Pasted image
20250301220243.png{]}{]} !{[}{[}Pasted image 20250301220302.png{]}{]}
\#\#\# Question 11: Write a shell script that reads numbers from the
user until a negative number is entered.

\begin{Shaded}
\begin{Highlighting}[]
\CommentTok{\#!/bin/bash}
\ControlFlowTok{while} \FunctionTok{true}\KeywordTok{;} \ControlFlowTok{do}
    \BuiltInTok{read} \AttributeTok{{-}p} \StringTok{"Enter a number: "} \VariableTok{num}
    \ControlFlowTok{if} \BuiltInTok{[} \VariableTok{$num} \OtherTok{{-}lt}\NormalTok{ 0 }\BuiltInTok{]}\KeywordTok{;} \ControlFlowTok{then}
        \ControlFlowTok{break}
    \ControlFlowTok{fi}
    \BuiltInTok{echo} \StringTok{"Square: }\VariableTok{$((num} \OperatorTok{*} \VariableTok{num))}\StringTok{"}
\ControlFlowTok{done}
\end{Highlighting}
\end{Shaded}

!{[}{[}Pasted image 20250301220454.png{]}{]}

\begin{center}\rule{0.5\linewidth}{0.5pt}\end{center}

\section{Part D - Common Interview
Questions}\label{part-d---common-interview-questions}

\begin{enumerate}
\def\labelenumi{\arabic{enumi}.}
\tightlist
\item
  What is an operating system, and what are its primary functions?
\item
  Explain the difference between process and thread.
\item
  What is virtual memory, and how does it work?
\item
  Describe the difference between multiprogramming, multitasking, and
  multiprocessing.
\item
  What is a file system, and what are its components?
\item
  What is a deadlock, and how can it be prevented?
\item
  Explain the difference between a kernel and a shell.
\item
  What is CPU scheduling, and why is it important?
\item
  How does a system call work?
\item
  What is the purpose of device drivers in an operating system?
\item
  Explain the role of the page table in virtual memory management.
\item
  What is thrashing, and how can it be avoided?
\item
  Describe the concept of a semaphore and its use in synchronization.
\item
  How does an operating system handle process synchronization?
\item
  What is the purpose of an interrupt in operating systems?
\item
  Explain the concept of a file descriptor.
\item
  How does a system recover from a system crash?
\item
  Describe the difference between a monolithic kernel and a microkernel.
\item
  What is the difference between internal and external fragmentation?
\item
  How does an operating system manage I/O operations?
\item
  Explain the difference between preemptive and non-preemptive
  scheduling.
\item
  What is round-robin scheduling, and how does it work?
\item
  Describe the priority scheduling algorithm. How is priority assigned
  to processes?
\item
  What is the shortest job next (SJN) scheduling algorithm, and when is
  it used?
\item
  Explain the concept of multilevel queue scheduling.
\item
  What is a process control block (PCB), and what information does it
  contain?
\item
  Describe the process state diagram and the transitions between
  different process states.
\item
  How does a process communicate with another process in an operating
  system?
\item
  What is process synchronization, and why is it important?
\item
  Explain the concept of a zombie process and how it is created.
\item
  Describe the difference between internal fragmentation and external
  fragmentation.
\item
  What is demand paging, and how does it improve memory management
  efficiency?
\item
  How does a memory management unit (MMU) work?
\item
  What is a system call, and how does it facilitate communication
  between user programs and the operating system?
\item
  Explain the concept of a race condition and how it can be prevented.
\item
  How does the fork () system call work in creating a new process in
  Unix-like operating systems?
\item
  How does process termination occur in Unix-like operating systems?
\item
  What is the role of the long-term scheduler in process scheduling?
\item
  Describe how a parent process can wait for a child process to finish
  execution.
\item
  What is the significance of the exit status of a child process in the
  wait () system call?
\end{enumerate}

\section{Part E - Scheduling and Process Management
Questions}\label{part-e---scheduling-and-process-management-questions}

\subsection{1. First-Come, First-Served (FCFS)
Scheduling}\label{first-come-first-served-fcfs-scheduling}

Consider the following processes with arrival times and burst times:

\begin{longtable}[]{@{}lll@{}}
\toprule\noalign{}
Process & Arrival Time & Burst Time \\
\midrule\noalign{}
\endhead
\bottomrule\noalign{}
\endlastfoot
P 1 & 0 & 5 \\
P 2 & 1 & 3 \\
P 3 & 2 & 6 \\
\end{longtable}

Calculate the average waiting time using First-Come, First-Served (FCFS)
scheduling. !{[}{[}Pasted image 20250301233236.png{]}{]} ---

\subsection{2. Shortest Job First (SJF)
Scheduling}\label{shortest-job-first-sjf-scheduling}

Consider the following processes with arrival times and burst times:

\begin{longtable}[]{@{}lll@{}}
\toprule\noalign{}
Process & Arrival Time & Burst Time \\
\midrule\noalign{}
\endhead
\bottomrule\noalign{}
\endlastfoot
P 1 & 0 & 3 \\
P 2 & 1 & 5 \\
P 3 & 2 & 1 \\
P 4 & 3 & 4 \\
\end{longtable}

Calculate the average turnaround time using Shortest Job First (SJF)
scheduling. !{[}{[}Pasted image 20250301233553.png{]}{]} ---

\subsection{3. Priority Scheduling}\label{priority-scheduling}

Consider the following processes with arrival times, burst times, and
priorities (lower number indicates higher priority):

\begin{longtable}[]{@{}llll@{}}
\toprule\noalign{}
Process & Arrival Time & Burst Time & Priority \\
\midrule\noalign{}
\endhead
\bottomrule\noalign{}
\endlastfoot
P 1 & 0 & 6 & 3 \\
P 2 & 1 & 4 & 1 \\
P 3 & 2 & 7 & 4 \\
P 4 & 3 & 2 & 2 \\
\end{longtable}

Calculate the average waiting time using Priority Scheduling.
!{[}{[}WhatsApp Image 2025-03-01 at 23.28.52\_bd738bb6 1.jpg{]}{]} ---

\subsection{4. Round Robin Scheduling}\label{round-robin-scheduling}

Consider the following processes with arrival times and burst times, and
the time quantum for Round Robin scheduling is 2 units:

\begin{longtable}[]{@{}lll@{}}
\toprule\noalign{}
Process & Arrival Time & Burst Time \\
\midrule\noalign{}
\endhead
\bottomrule\noalign{}
\endlastfoot
P 1 & 0 & 4 \\
P 2 & 1 & 5 \\
P 3 & 2 & 2 \\
P 4 & 3 & 3 \\
\end{longtable}

Calculate the average turnaround time using Round Robin scheduling.
!{[}{[}Pasted image 20250301233738.png{]}{]} ---

\subsection{5. Fork System Call
Scenario}\label{fork-system-call-scenario}

Consider a program that uses the \texttt{fork()} system call to create a
child process. Initially, the parent process has a variable \texttt{x}
with a value of \texttt{5}. After forking, both the parent and child
processes increment the value of \texttt{x} by \texttt{1}.

What will be the final values of \texttt{x} in the parent and child
processes after the \texttt{fork()} call? !{[}{[}Pasted image
20250301234027.png{]}{]} !{[}{[}WhatsApp Image 2025-03-01 at
23.28.53\_08256142.jpg{]}{]}
